\documentclass[11pt,letterpaper]{article}
\usepackage[margin=1in]{geometry}
\usepackage{amsmath,amssymb,mathtools}
\usepackage[T1]{fontenc}
\usepackage[utf8]{inputenc}
\usepackage{lmodern}
\DeclareUnicodeCharacter{0398}{\ensuremath{\Theta}}
\DeclareUnicodeCharacter{03B3}{\ensuremath{\gamma}}
\DeclareUnicodeCharacter{03B5}{\ensuremath{\varepsilon}}
\DeclareUnicodeCharacter{2113}{\ensuremath{\ell}}
\DeclareUnicodeCharacter{00B2}{\ensuremath{{}^2}}
\DeclareUnicodeCharacter{00B3}{\ensuremath{{}^3}}
\usepackage{microtype}
\usepackage{booktabs,longtable,array}
\usepackage{enumitem}
\setlist{itemsep=3pt,topsep=5pt}
\usepackage{xurl}
\usepackage[hidelinks,unicode]{hyperref}
\hypersetup{pdftitle={Asymptotic amplitude and all orders for relaxed binary trees},pdfauthor={},pdfsubject={A082161 spectral asymptotics and inverse thresholds}}
\usepackage{bookmark}
\setlength{\parskip}{4pt}
\setlength{\parindent}{0pt}
\setlength{\emergencystretch}{2em}
\allowdisplaybreaks[2]
\providecommand{\tightlist}{\setlength{\itemsep}{0pt}\setlength{\parskip}{0pt}}
\newcommand{\passthrough}[1]{#1}
\title{Asymptotic amplitude and all orders\\for relaxed binary trees}
\author{}
\date{2 October 2026}
\begin{document}
\maketitle
\vspace{-1.5em}
\begin{abstract}
The number of relaxed binary trees is known to have order
$n!4^n\exp(3zn^{1/3})n$, where $z$ is the largest negative Airy zero.
We give a spectral argument for convergence to a positive limiting amplitude and for a full expansion in powers of $n^{-1/3}$.
An exact time-dependent gauge turns the counting recurrence into a symmetric Jacobi operator followed by a diagonal contraction.
A finite Airy quasimode, the continuum spectral gap, and summable ground-mode drift establish norm tracking.
A Catalan kernel estimate recovers the boundary coordinate.
Polynomial--Airy approximate solutions and a zero-limit projection bootstrap then transfer every finite formal order to the exact sequence.
The known lower bound is used to prove strict positivity of the amplitude.
We also derive Lambert $W$ inverse expansions and distinguish them from certified discrete threshold decisions.
The accompanying reproducibility archive checks the exact algebra and numerical evidence; its amplitude estimates are not rigorous enclosures.
\end{abstract}
\tableofcontents
\newpage
\section{Statement and scope}\label{statement-and-scope}

Let \(a_n\) be OEIS A082161, the number of relaxed binary trees with \(n\) internal nodes. Its defining triangular recurrence is

\[
r_{n,m}=r_{n,m-1}+(m+1)r_{n-1,m},\qquad
r_{n,0}=1,\quad r_{n,m}=0\ (m>n),\quad a_n=r_{n,n}.
\]

Let \(z=a_1^{\operatorname{Ai}}=-2.338107410459767\ldots\) be the largest real zero of \(\operatorname{Ai}\). (This Airy zero is denoted by \(z\) throughout, to avoid confusing it with the counting sequence.) Set

\[
a=2^{-1/3}z,\qquad F(x)=\operatorname{Ai}(z+2^{1/3}x).
\]

\textbf{Theorem.} There is a constant \(\gamma >0\) such that, for each integer \(M\ge 0\),

\[
a_n=\gamma n!4^n e^{3zn^{1/3}}n
\left(1+\sum_{k=1}^M c_kn^{-k/3}
+O_M(n^{-(M+1)/3})\right).
\tag{F1}
\]

The coefficients belong to \(\mathbb Q[z]\) and are determined by the polynomial--Airy recursion in §8. The first logarithmic coefficients are

\[
\log\frac{a_n}{\gamma n!4^n e^{3zn^{1/3}}n}
=\frac{53z^2}{90}n^{-1/3}
+\frac{44z}{27}n^{-2/3}
+\left(\frac{141}{140}-\frac{1304z^3}{42525}\right)n^{-1}
+O(n^{-4/3}).
\tag{F2}
\]

Consequently

\[
\begin{aligned}
c_1&=53z^2/90,\\
c_2&=z(2809z^3+26400)/16200,\\
c_3&=(1042139z^6+28444320z^3+30836700)/30618000.
\end{aligned}
\tag{F3}
\]

Numerically these are \(3.2193061325653844\), \(1.3717168742447035\), and \(-5.306490818002809\). No rigorous numerical enclosure or elementary evaluation of \(\gamma\) is asserted.

The counting recurrence, the Θ estimate, and the conjectural limiting amplitude and all-orders shape are due to Andrew Elvey Price, Wenjie Fang, and Michael Wallner, \href{https://arxiv.org/pdf/1908.11181}{Compacted binary trees admit a stretched exponential}, especially Theorem 1.1 and §3.4. Their paper explicitly leaves existence of the amplitude open and reports the fit \(166.95208957\). The contribution here is a convergence/stability argument for the relaxed recurrence. It does not extend without further work to compacted trees, whose recurrence is different. It proves a Poincaré expansion to every finite algebraic order, not an exponentially improved transseries or a canonical interpolation between integers.

\section{Exact recurrence and contractive symmetrization}\label{exact-recurrence-and-contractive-symmetrization}

Put

\[
d_{N,j}=\frac{r_{(N+j)/2,(N-j)/2}}{((N+j)/2)!}
\]

when \(0\le j\le N\) and \(j\equiv N\pmod 2\), and zero otherwise. Direct substitution gives

\[
d_{0,j}=\mathbf 1_{j=0},\qquad d_{N,-1}=0,
\qquad
d_{N,j}=\frac{N-j+2}{N+j}d_{N-1,j-1}+d_{N-1,j+1}.
\tag{F4}
\]

In particular \(a_n=n!d_{2n,0}\). Values outside the admissible support are interpreted as zero.

For \(N\ge 1\), define on \(0\le j\le N+1\)

\[
D_N(j)^2=\frac{\Gamma(N+2)\Gamma(N+1)}
{\Gamma(N-j+2)\Gamma(N+j+1)},\qquad D_N(0)=1,
\]

and \(v_N(j)=d_{N,j}/D_N(j)\). All vectors are padded with zeros in \(\ell ^2(\mathbb Z_{\ge 0})\). Define \(S_N\) to be the symmetric Jacobi matrix on \(0,\ldots ,N+1\), padded by zero, with edge weights

\[
b_{N,j}=\sqrt{\frac{N-j+2}{N+j}},\quad 1\le j\le N+1,
\]

and define a diagonal contraction by

\[
Q_N(j)=\sqrt{1-\frac{j(j-1)}{N(N+1)}}
\quad(0\le j\le N+1),\qquad Q_N(j)=0\quad(j>N+1).
\]

The factorial identities

\[
\frac{D_N(j)}{D_N(j-1)}=b_{N,j},\qquad
\frac{D_{N-1}(j)}{D_N(j)}=Q_N(j)
\]

on the support where the second ratio is needed establish the exact identity

\[
v_N=S_NQ_Nv_{N-1}.
\tag{F5}
\]

Both \(b\) and \(Q\) lie in \([0,1]\). If \(J\) denotes the half-line adjacency matrix, then \(0\le S_NQ_N\le J\) entrywise. If \(\lambda _N\) is the largest eigenvalue of \(S_N\), then \(\Vert S_NQ_N\Vert \le \lambda _N\). The finite Jacobi matrix is irreducible, so its top eigenvector \(\psi _N\) can be taken positive and of norm one. Its spectrum is symmetric about zero because it exchanges the two parity subspaces. In particular its operator norm is \(\lambda _N\).

\section{The coarse Airy limit and a uniform spectral gap}\label{the-coarse-airy-limit-and-a-uniform-spectral-gap}

Write \(\varepsilon =N^{-1/3}\) and \(x_j=(j+1)\varepsilon\). For vectors supported on the finite interval \(0,\ldots ,N+1\), the quadratic form is exactly

\[
\begin{aligned}
\langle v,(2I-S_N)v\rangle={}&
\sum_{j=0}^{N}b_{N,j+1}|v_{j+1}-v_j|^2+|v_0|^2\\
&+\sum_{j=1}^{N}(2-b_{N,j}-b_{N,j+1})|v_j|^2
+(2-b_{N,N+1})|v_{N+1}|^2.
\end{aligned}
\tag{F6}
\]

Here \(|v_0|^2\) is the first Dirichlet-gradient interval, with fictitious value \(v_{-1}=0\); no nonphysical edge coefficient \(b_{N,0}\) is introduced.

For \(1\le j\le N+1\),

\[
1-b_{N,j}
=\frac{2(j-1)}{(N+j)(1+b_{N,j})}
\ge\frac{j-1}{N+j}.
\]

It follows from (F6) that, for universal positive constants \(c,C\),

\[
\varepsilon^{-2}\langle v,(2I-S_N)v\rangle
\ge c\sum_{j=0}^{N+1}x_j|v_j|^2-C\varepsilon\|v\|^2.
\tag{F7}
\]

On each fixed \(x\) window the edge weights tend uniformly to one, and

\[
\varepsilon^{-2}(2-b_{N,j}-b_{N,j+1})\longrightarrow 2x_j
\]

uniformly there. Interpolate linearly with value \(\varepsilon ^{-1/2}v_j\) at \(x_j\) and zero at zero. Add one zero node immediately after the final finite node and extend by zero thereafter; the endpoint term in (F6) controls this final interpolation interval. The gradient terms in (F6) give local \(H^1\) bounds for any sequence with bounded scaled energy. Equation (F7) gives uniformly small \(L^2\) tails outside a large fixed window. Local compactness, followed by this tail bound, gives strong \(L^2\) compactness. The boundary interval enforces the trace zero at zero. The resulting liminf form is

\[
\int_0^\infty (|f'(x)|^2+2x|f(x)|^2)\,dx.
\]

For completeness, convergence of norms in this interpolation follows first on a fixed window from the local gradient bound and the Riemann-sum identity, and then globally by (F7). Sampling smooth compactly supported functions with zero boundary value gives the reverse, limsup bound. Applying the elementary min--max principle in both directions therefore gives convergence of each fixed low eigenvalue of \(\varepsilon ^{-2}(2I-S_N)\) to that of

\[
\mathcal H=-\frac{d^2}{dx^2}+2x,
\qquad f(0)=0.
\]

The latter eigenvalues are \(\mu _k=-2^{2/3}z_k\), where \(z_1=z>z_2>\ldots\) are the negative Airy zeros. Thus, if \(\lambda _{2,N}\) is the second-largest positive eigenvalue for large \(N\),

\[
\lambda_N=2+2a\varepsilon^2+o(\varepsilon^2),\qquad
\lambda_N-\lambda_{2,N}
\sim 2^{2/3}(z-z_2)\varepsilon^2.
\tag{F8}
\]

As a map between opposite parity subspaces, \(S_N\) has singular values equal to its positive eigenvalues (and possibly additional zeros). Hence, after deleting its top singular direction, its norm divided by \(\lambda _N\) is at most \(1-cN^{-2/3}\) for all sufficiently large \(N\). This parity formulation is essential: the negative eigenvalue \(-\lambda _N\) is not a second slow direction within the occupied parity space.

\section{Quantitative estimates from a finite quasimode}\label{quantitative-estimates-from-a-finite-quasimode}

The Airy equation reads \(F''=2(x+a)F\). Define

\[
\begin{aligned}
F_2(x)&=\frac{2(-a+2x)}{15}F(x)
+\frac{x(2a+x)}{15}F'(x),\\
F_3(x)&=\frac{F(x)-xF'(x)}{3},\\
f_N(x)&=F(x)+\varepsilon^2F_2(x)+\varepsilon^3F_3(x),\\
\widehat\lambda_N&=2+2a\varepsilon^2+3\varepsilon^3
+\frac{13a^2}{15}\varepsilon^4+\frac{5a}{3}\varepsilon^5.
\end{aligned}
\tag{F9}
\]

Choose a fixed smooth cutoff \(\chi\), equal to one on \([0,1]\) and zero on \([2,\infty )\), and sample \(\chi (j/N^{1/2})f_N(x_j)\). Let \(\widehat f_N\) denote its \(\ell ^2\) normalization. Changing or omitting the cutoff in local calculations produces an error smaller than every inverse power of \(N\).

Here is an explicit way to check and bound the quasimode residual. At a sample point the two edge coefficients are

\[
b_-(x,\varepsilon)=\sqrt{\frac{1-x\varepsilon^2+3\varepsilon^3}
{1+x\varepsilon^2-\varepsilon^3}},\qquad
b_+(x,\varepsilon)=\sqrt{\frac{1-x\varepsilon^2+2\varepsilon^3}
{1+x\varepsilon^2}}.
\]

Expand \(b_- f_N(x-\varepsilon )+b_+ f_N(x+\varepsilon )-\widehat\lambda _N f_N(x)\) through degree five in \(\varepsilon\), reducing every derivative using \(F''=2(x+a)F\). Every coefficient vanishes. At the left boundary the omitted sample is zero because \(F(0)=F_2(0)=F_3(0)=0\). On the support of the cutoff, \(x\varepsilon ^2=O(N^{-1/2})\); Taylor remainders are bounded by \(C\varepsilon ^6\) times a fixed polynomial in \(x\) multiplying Airy decay. All such polynomial--Airy envelopes are square-integrable. Derivatives of the cutoff contribute only \(\exp(-cN^{1/4})\) times a power of \(N\). Since the unnormalized sampled norm is asymptotic to \(\varepsilon ^{-1/2}\Vert F\Vert _{L^2}\), this proves

\[
\|(S_N-\widehat\lambda_N)\widehat f_N\|=O(\varepsilon^6).
\tag{F10}
\]

The Airy gap (F8) identifies the nearby eigenvalue with \(\lambda _N\); projecting the residual on the complement of \(\psi _N\) then gives

\[
\lambda_N=\widehat\lambda_N+O(\varepsilon^6),\qquad
\|\psi_N-\widehat f_N\|=O(\varepsilon^4).
\tag{F11}
\]

The sign is fixed by the positive leading Airy profile. Several useful consequences follow without a separate quantitative localization theorem:

\[
\begin{aligned}
\|\psi_N-\psi_{N-1}\|&=O(N^{-1}),\\
\|(I-Q_N)\psi_N\|&=O(N^{-4/3}),\\
\psi_N(0)&=\sqrt2\,N^{-1/2}(1+O(N^{-1/3})).
\end{aligned}
\tag{F12}
\]

For the first, the explicitly normalized samples in (F9) have \(\ell ^2\) derivative \(O(N^{-1})\) with respect to continuous \(N\); the two approximation errors in (F11) are smaller. For the second, use \(1-\sqrt{1-t}\le t\) for \(0\le t\le 1\), the finite polynomial moments of the sampled Airy profile, and (F11). For the endpoint, the pointwise error is at most the \(\ell ^2\) error, and

\[
I:=\int_0^\infty F(x)^2\,dx=2^{-1/3}\operatorname{Ai}'(z)^2,
\qquad \frac{F'(0)}{\sqrt I}=\sqrt2.
\]

A Riemann-sum estimate with an integrable derivative gives a relative \(O(\varepsilon )\) normalization error; the endpoint Taylor error is \(O(\varepsilon ^2)\) and the normalized quasimode error divided by the leading endpoint is \(O(\varepsilon ^{5/2})\). These identities therefore verify the quantitative endpoint estimate in (F12). Standard Airy decay needed above follows, for example, directly from the decaying solution of the Airy ODE; only its polynomially weighted integrability and exponential tail are used.

\section{Convergence of the scalar amplitude}\label{convergence-of-the-scalar-amplitude}

Set

\[
P_N=\prod_{i=1}^N\lambda_i,\qquad
u_N=v_N/P_N,\qquad M_N=S_NQ_N/\lambda_N.
\]

Then \(u_N=M_Nu_{N-1}\), \(\Vert M_N\Vert \le 1\), and \(\Vert u_N\Vert \le 1\). Let \(E_p\) be projection onto height parity \(p\). The positive eigenvector has equal squared mass on the two parities: this follows by taking the scalar product of each half of \(S_N\psi _N=\lambda _N\psi _N\) with that half. Therefore

\[
g_N=\sqrt2 E_{N\bmod2}\psi_N,\qquad
h_N=\sqrt2 E_{(N-1)\bmod2}\psi_N
\]

are unit vectors, with \(S_Nh_N=\lambda _Ng_N\). By (F12),

\[
\|h_N-g_{N-1}\|=O(N^{-1}),\qquad
\|(I-Q_N)h_N\|+\|(I-Q_N)g_{N-1}\|=O(N^{-4/3}).
\tag{F13}
\]

The second estimate for \(g_{N-1}\) uses the first estimate only for its approximation by explicit Airy samples, rather than the crude \(O(N^{-1})\) difference: direct use of (F11) at \(N-1\) gives the stated \(O(N^{-4/3})\) estimate with \(Q_N\). It also follows by observing that \(Q_N\ge Q_{N-1}\) pointwise on the old support and treating its negligible eigenvector tail.

Decompose

\[
u_N=A_Ng_N+w_N,\qquad w_N\perp g_N.
\]

Write \(\Pi _N=I-|g_N\rangle \langle g_N|\) on the occupied parity subspace. The singular gap from §3 gives, for every vector \(y\) on the preceding parity,

\[
\|\Pi_NM_Ny\|
\le (1-cN^{-2/3})\|Q_Ny\|
\le (1-cN^{-2/3})\|y\|.
\tag{F14}
\]

Indeed \(\Pi _NS_N=S_N(I-|h_N\rangle \langle h_N|)\) on that parity; this projection removes the top singular direction, regardless of whether \(y\) is perpendicular to \(h_N\). Also, (F13) implies

\[
\|M_Ng_{N-1}-g_N\|=O(N^{-1}).
\]

Applying (F14) to \(w_{N-1}\) and the latter estimate to the scalar term yields

\[
\|w_N\|\le (1-cN^{-2/3})\|w_{N-1}\|+CN^{-1}.
\tag{F15}
\]

Iteration gives \(\Vert w_N\Vert =O(N^{-1/3})\). One general estimate used here and later is

\[
\sum_{k=N_0}^N k^{-r}
\prod_{i=k+1}^N(1-ci^{-2/3})=O_r(N^{2/3-r}).
\tag{F16}
\]

To prove it, split at \(N/2\). The earlier half is exponentially small in \(N^{1/3}\) times a polynomial; on the later half bound the product by \(\exp(-c'(N-k)N^{-2/3})\) and sum a geometric series. A finite initial segment makes no difference.

The scalar equation is

\[
A_N=t_NA_{N-1}+\langle k_N,w_{N-1}\rangle,
\quad
t_N=\langle Q_Nh_N,g_{N-1}\rangle,
\quad k_N=Q_Nh_N-g_{N-1}.
\tag{F17}
\]

The unit-vector identity \(\langle h,g\rangle =1-\Vert h-g\Vert ^2/2\) and (F13) show

\[
t_N=1+O(N^{-4/3}),\qquad \|k_N\|=O(N^{-1}).
\tag{F18}
\]

Since \(|A_N|\le 1\), (F15), (F17), and (F18) give

\[
A_N-A_{N-1}=O(N^{-4/3}).
\]

Thus

\[
A_N=A_\infty+O(N^{-1/3}).
\tag{F19}
\]

Positivity gives \(A_N\ge 0\) and hence \(A_\infty \ge 0\); strict positivity is established in §7.

\section{A boundary smoothing estimate}\label{a-boundary-smoothing-estimate}

Norm tracking alone cannot read a coordinate whose leading size is \(N^{-1/2}\). A short final block of the original positive kernel provides the missing smoothing.

Let \(U(N,k)=M_N\cdots M_{k+1}\). For \(0\le \ell \le N^{2/3}\), positivity and \(S_iQ_i\le J\) give

\[
\|e_0^TU(N,N-\ell)\|_2
\le\frac{\|e_0^TJ^\ell\|_2}
{\lambda_{N-\ell+1}\cdots\lambda_N}.
\]

The squared numerator is the number of nonnegative walks of length \(2\ell\) returning to zero:

\[
\|e_0^TJ^\ell\|_2^2=(J^{2\ell})_{0,0}
=\frac{1}{\ell+1}\binom{2\ell}{\ell}
\le C4^\ell(\ell+1)^{-3/2}.
\]

Equation (F11) implies

\[
\frac{2^\ell}{\lambda_{N-\ell+1}\cdots\lambda_N}
\le\exp(C\ell N^{-2/3})\le C'
\]

over the same range. Therefore

\[
\|e_0^TU(N,N-\ell)\|_2\le C(\ell+1)^{-3/4}.
\tag{F20}
\]

The exact error recurrence for \(w_N\) is

\[
w_N=M_Nw_{N-1}+f_N^{\rm err},\qquad
f_N^{\rm err}=A_{N-1}M_Ng_{N-1}-A_Ng_N,
\quad \|f_N^{\rm err}\|=O(N^{-1}).
\]

Use Duhamel over \(L=\lfloor N^{2/3}\rfloor\) steps. Since all times in this block are comparable with \(N\), (F15) and (F20) imply

\[
\begin{aligned}
|w_N(0)|
&\le CL^{-3/4}\|w_{N-L}\|
+CN^{-1}\sum_{\ell=0}^{L-1}(\ell+1)^{-3/4}\\
&=O(N^{-5/6}).
\end{aligned}
\tag{F21}
\]

For even \(N\), (F12) gives \(g_N(0)=2N^{-1/2}(1+O(N^{-1/3}))\).

\section{The limiting constant and its strict positivity}\label{the-limiting-constant-and-its-strict-positivity}

From (F11),

\[
\log\lambda_N=\log2+aN^{-2/3}+\frac32N^{-1}+O(N^{-4/3}).
\]

Summation gives

\[
P_N=\kappa\,2^N e^{3aN^{1/3}}N^{3/2}(1+O(N^{-1/3})),
\qquad \kappa>0.
\tag{F22}
\]

If desired, a convergent expression is

\[
\log\kappa=a\zeta(2/3)+\tfrac32\gamma_E
+\sum_{N=1}^\infty\left[
\log(\lambda_N/2)-aN^{-2/3}-\frac{3}{2N}\right].
\tag{F23}
\]

Here \(\gamma _E\) is Euler's constant, not the counting amplitude. This is an existence formula; it does not by itself give efficient certified numerical evaluation.

Equations (F19)--(F22), for even \(N\), now give

\[
d_{N,0}=2\kappa A_\infty\,2^N e^{3aN^{1/3}}N
+O\left(2^N e^{3aN^{1/3}}N^{2/3}\right).
\tag{F24}
\]

The known lower half of the Elvey Price--Fang--Wallner Θ theorem states that \(d_{2n,0}\) is bounded below by a positive constant times \(4^n e^{3zn^{1/3}}n\) for all sufficiently large \(n\). If \(A_\infty =0\), (F24) contradicts that bound. Hence \(A_\infty >0\) and

\[
\gamma=4\kappa A_\infty>0.
\tag{F25}
\]

Taking \(N=2n\) in (F24) proves (F1) with \(M=0\). This completes the amplitude-convergence argument, using the explicit prior lower-bound theorem and the elementary Airy facts used above.

\section{Formal profiles at every finite order}\label{formal-profiles-at-every-finite-order}

Work now with the original recurrence (F4), rather than the symmetric gauge. Seek

\[
Z_{N,j}=H_N\sum_{k=0}^K\varepsilon^kG_k(x_j),\qquad
G_0=F,\quad G_k(0)=G_k'(0)=0\ (k\ge1),
\]

with

\[
\frac{H_N}{H_{N-1}}
\sim2\left(1+\sum_{m\ge2}s_m\varepsilon^m\right).
\tag{F26}
\]

At the previous time the two arguments and powers are

\[
y_\pm=(x\pm\varepsilon)(1-\varepsilon^3)^{-1/3},
\qquad
\varepsilon_{N-1}^k=\varepsilon^k(1-\varepsilon^3)^{-k/3},
\]

and the downward coefficient is

\[
\frac{1-x\varepsilon^2+3\varepsilon^3}
{1+x\varepsilon^2-\varepsilon^3}.
\]

After the leading orders, the coefficient at degree \(m\) has the form

\[
\mathcal LG_{m-2}-2s_mF+R_m(x)F+S_m(x)F',
\quad
\mathcal L=\frac{d^2}{dx^2}-2(x+a),
\tag{F27}
\]

where \(R_m,S_m\) are known polynomials over \(\mathbb Q[a]\), formed from earlier orders. In particular \(s_2=a\). This form follows by Taylor expansion: the new profile first appears through the second-order spatial operator; its lower, zeroth-order terms cancel between the two sides of (F4). Time-change terms involving this new profile enter only at later degree.

To solve \(\mathcal L(PF+QF')=RF+SF'\), compute

\[
\mathcal L(PF+QF')
=(P''+4(x+a)Q'+2Q)F+(2P'+Q'')F'.
\]

Eliminating \(P'=(S-Q'')/2\) reduces the problem to

\[
-\tfrac12 Q'''+4(x+a)Q'+2Q=R-\tfrac12S'.
\tag{F28}
\]

On polynomials of degree at most \(d\), the operator on the left is upper triangular by degree with nonzero diagonal entries \(4j+2\), \(0\le j\le d\); hence it is invertible. The new scalar \(s_m\) adds \(2s_m\) to the forcing \(R\), thus adds exactly \(s_m\) to the constant term of \(Q\). Choose it uniquely to enforce \(Q(0)=0\), which is the boundary condition. Integrating \(P'\) leaves one free constant; choose it so \(P(0)+Q'(0)=0\), which enforces the derivative normalization. This proves existence and uniqueness at every formal order within the stated normalization, not merely consistency of finitely many guessed coefficients.

The first scalar coefficients are

\[
s_2=a,\quad s_3=\frac43,\quad s_4=\frac{29a^2}{270},\quad
s_5=-\frac{23a}{81},\quad
s_6=\frac{74385-8783a^3}{85050}.
\tag{F29}
\]

For example,

\[
G_1=-\tfrac23x^2F,\quad
G_2=\frac{2(2a+15x^4+111x)}{135}F
+\frac{x(-8a+51x)}{270}F'.
\]

The amplitude factor can be chosen as a genuine smooth function of \(N\) with finite logarithmic expansion

\[
\log H_N=N\log2+3aN^{1/3}+\frac43\log N
+\sum_{k=1}^J h_kN^{-k/3}.
\tag{F30}
\]

Expanding its exact one-step ratio matches any prescribed finite number of coefficients in (F26). At degree \(\varepsilon ^{k+3}\), the new \(h_k\) has coefficient \(-k/3\), so the construction is again triangular and nonsingular. Crucially the leading multiplicative constant in (F30) is fixed to one for every truncation. This avoids an artificial truncation-dependent amplitude.

\section{Approximate solutions in the normalized norm}\label{approximate-solutions-in-the-normalized-norm}

Fix any \(p>1\). Take enough orders in §8, multiply the sampled profile by the cutoff \(\chi (j/N^{1/2})\), and restrict to the occupied parity. Call the result \(Z_N\). Choose the truncation so that its local formal residual starts at a power \(N^{-p'}\) with \(p'\ge p\). Define

\[
z_N=D_N^{-1}Z_N/P_N.
\]

Then

\[
\|z_N\|=O(1),\qquad
r_N:=z_N-M_Nz_{N-1}=O(N^{-p})\quad\hbox{in }\ell^2.
\tag{F31}
\]

Here are the needed global bounds. On \(j\le 2N^{1/2}\), the product formula for \(D_N\) gives \(D_N^{-1}\le C\) uniformly: \(-\log D_N(j)\le Cj^2/N\le C'\). On any fixed scaled window \(D_N^{-1}=1+O(N^{-1/3}(1+x)^2)\). Taylor remainders in the rational coefficients and the time dilation are bounded by \(N^{-p}\) times a fixed polynomial--Airy envelope on the cutoff support. Cutoff errors decay like \(\exp(-cN^{1/4})\) times a polynomial, also after multiplication by \(D_N^{-1}\). Finally,

\[
H_N/P_N\sim\kappa^{-1}N^{-1/6},
\]

while a parity-restricted sampled profile has norm asymptotic to \(N^{1/6}\sqrt{I/2}\). These bounds prove (F31) by direct summation. The same reasoning proves

\[
\langle g_N,z_N\rangle\longrightarrow
q_0:=\kappa^{-1}\sqrt{I/2}>0.
\tag{F32}
\]

Indeed, after the natural sample normalization, both profiles converge to the same positive Airy function. Neither (F31) nor (F32) needs a pointwise estimate on the unknown solution.

\section{A scalar bootstrap to every finite order}\label{a-scalar-bootstrap-to-every-finite-order}

Set

\[
C=A_\infty/q_0,\qquad e_N=u_N-Cz_N.
\]

Then \(e_N=M_Ne_{N-1}-Cr_N\) and the ground projection tends to zero. Decompose

\[
e_N=\alpha_Ng_N+W_N,\qquad W_N\perp g_N,
\qquad \alpha_N\longrightarrow0.
\]

The exact same estimates (F14), (F17), and (F18) now give

\[
\begin{aligned}
\|W_N\|&\le(1-cN^{-2/3})\|W_{N-1}\|
+CN^{-1}|\alpha_{N-1}|+CN^{-p},\\
|\alpha_N-\alpha_{N-1}|&\le
CN^{-4/3}|\alpha_{N-1}|+CN^{-1}\|W_{N-1}\|+CN^{-p}.
\end{aligned}
\tag{F33}
\]

Suppose initially that \(\alpha _N=O(N^{-q})\), where \(q\ge 0\). Formula (F16) gives

\[
\|W_N\|=O(N^{-q-1/3}+N^{2/3-p}).
\]

Substitute this into the second inequality of (F33). Its right-hand side is summable because \(p>1\); summing from \(N+1\) to infinity and using the zero limiting projection gives

\[
|\alpha_N|=O(N^{-q-1/3}+N^{1-p}).
\tag{F34}
\]

The smaller tail \(O(N^{2/3-p})\) has been absorbed by \(O(N^{1-p})\). Start with \(q=0\), which holds since \(u_N,z_N\) are bounded. Finitely many repetitions of (F34) reach

\[
\alpha_N=O(N^{1-p}),\qquad
\|W_N\|=O(N^{2/3-p}),\qquad
\|e_N\|=O(N^{1-p}).
\tag{F35}
\]

There is no boundary-at-infinity assumption beyond the already proved, explicit limit (F32) and the choice of \(C\). No inverse operator with a small gap is silently assumed uniformly bounded.

Apply Duhamel directly to the inhomogeneous \(e\) recurrence on the final block \(L=\lfloor N^{2/3}\rfloor\). By (F20), (F31), and (F35),

\[
\begin{aligned}
|e_N(0)|&\le CL^{-3/4}N^{1-p}
+CN^{-p}\sum_{\ell=0}^{L-1}(\ell+1)^{-3/4}\\
&=O(N^{1/2-p})+O(N^{1/6-p})
=O(N^{1/2-p}).
\end{aligned}
\tag{F36}
\]

The main endpoint is of size \(N^{-1/2}\), so the relative endpoint error is \(O(N^{1-p})\). For the remainder at order \(M\), take \(p\ge 1+(M+1)/3\) and enough formal terms. Since \(p\) can be taken as large as desired, every prescribed finite algebraic order of the explicit approximation transfers to the exact sequence. All truncations have the same \(q_0\), hence the same \(C\) and the same \(\gamma =2CF'(0)=4\kappa A_\infty\). This proves (F1).

\section{Coefficient extraction and computational checks}\label{coefficient-extraction-and-computational-checks}

Solving the scalar ratio matching (F30) using (F29) gives

\[
h_1=53a^2/45,\qquad h_2=79a/27,\qquad
h_3=143/210-5216a^3/42525.
\]

The normalized endpoint profile is

\[
\frac{\sum_k\varepsilon^kG_k(\varepsilon)}{F'(0)\varepsilon}
=1+\frac a3\varepsilon^2+\frac43\varepsilon^3+O(\varepsilon^4).
\]

Consequently the logarithmic endpoint corrections in time \(N\) are

\[
\frac{53a^2}{45}N^{-1/3}
+\frac{88a}{27}N^{-2/3}
+\left(\frac{141}{70}-\frac{5216a^3}{42525}\right)N^{-1}.
\]

Substituting \(N=2n\) and \(a=2^{-1/3}z\) yields (F2), and exponentiation gives (F3).

\section{The rational coefficient field}\label{the-rational-coefficient-field}

The recursion immediately puts every coefficient at time \(N\) in \(\mathbb Q[a]\), where \(a=2^{-1/3}z\). The substitution \(N=2n\) initially gives the weaker field \(\mathbb Q(2^{1/3})[z]\). To justify the stronger field \(\mathbb Q[z]\) asserted in the theorem, one needs the following mod-3 grading.

Let \(\phi_a\) be the formal power-series solution of

\[
\phi_a''=2(x+a)\phi_a,\qquad \phi_a(0)=0,\quad\phi_a'(0)=1.
\]

Every coefficient of \(\phi_a\) belongs to \(\mathbb Q[a]\). At the special Airy parameter used in the proof, this is exactly \(F(x)/F'(0)\). For any primitive cube root of unity \(\omega\), uniqueness of the initial-value solution gives

\[
\phi_{\omega a}(\omega x)=\omega\phi_a(x),\qquad
\phi'_{\omega a}(\omega x)=\phi'_a(x).
\tag{G1}
\]

Indeed, the left side of the first identity, as a function of \(x\), obeys the same ODE as \(\phi_a\) because \(\omega^3=1\), and its derivative at zero is \(\omega\).

Divide every formal profile in §8 by \(F'(0)\) and denote the resulting formal sum by

\[
\Phi(a,\varepsilon,x)=\sum_{k\ge0}\varepsilon^k\mathcal G_k(a,x),
\qquad \mathcal G_0=\phi_a.
\]

Its normalization is \(\Phi(a,\varepsilon,0)=0\) and \(\partial_x \Phi(a,\varepsilon,0)=1\). Denote the half-ratio of the amplitude factors by

\[
R(a,\varepsilon)=1+\sum_{m\ge2}s_m(a)\varepsilon^m.
\]

In the formal recurrence, the rational coefficient depends only on \(x \varepsilon^2\) and \(\varepsilon^3\); the dilation depends only on \(\varepsilon^3\); and the shifted arguments are \((x\pm \varepsilon)(1-\varepsilon^3)^{-1/3}\). Consequently the transformation

\[
(a,\varepsilon,x)\longmapsto(\omega a,\omega\varepsilon,\omega x)
\]

preserves the recurrence. The transformed profile

\[
\widetilde\Phi(a,\varepsilon,x)
=\omega^{-1}\Phi(\omega a,\omega\varepsilon,\omega x)
\]

has the same leading solution and the same boundary and derivative normalizations, by (G1). The order-by-order polynomial recursion in §8 is unique under exactly these normalizations. Therefore

\[
R(\omega a,\omega\varepsilon)=R(a,\varepsilon),\qquad
\Phi(\omega a,\omega\varepsilon,\omega x)
=\omega\Phi(a,\varepsilon,x).
\tag{G2}
\]

In particular,

\[
s_m(\omega a)=\omega^{-m}s_m(a).
\tag{G3}
\]

The uniquely determined coefficients \(h_k(a)\) of the finite logarithmic amplitude satisfy the same grading,

\[
h_k(\omega a)=\omega^{-k}h_k(a).
\tag{G4}
\]

One way to see (G4) without manipulating a logarithm branch is to use its ratio equation. The leading ratio contributions involve \(3a/\varepsilon\), multiplied by a series in \(\varepsilon^3\), and \(\log(1-\varepsilon^3)\). They are invariant under the transformation. The term involving \(h_k \varepsilon^k\) first enters at degree \(k+3\) with nonzero coefficient \(-k/3\), so the unique recursive solution must have the indicated grading.

The normalized endpoint factor is

\[
E(a,\varepsilon)=\frac{\Phi(a,\varepsilon,\varepsilon)}{\varepsilon}.
\]

Equation (G2) shows \(E(\omega a,\omega \varepsilon)=E(a,\varepsilon)\). The constant term of \(E\) is one. Thus every coefficient of both \(E\) and \(\log E\) has the same mod-3 grading. Combining it with (G4), the time-\(N\) logarithmic and relative correction coefficients \(b_k(a)\) obey

\[
b_k(\omega a)=\omega^{-k}b_k(a).
\tag{G5}
\]

Since \(b_k\) is a polynomial with rational coefficients, every nonzero monomial \(a^r\) in it satisfies

\[
r+k\equiv0\pmod3.
\]

Finally its contribution after changing variables is

\[
a^rN^{-k/3}
=2^{-(r+k)/3}z^r n^{-k/3}.
\]

The exponent \((r+k)/3\) is an integer, so the factor of two is rational. This proves that every forward correction coefficient at time \(n\) lies in \(\mathbb Q[z]\), rather than merely in \(\mathbb Q(2^{1/3})[z]\).

\section{Asymptotic inversion and discrete thresholds}\label{asymptotic-inversion-and-discrete-thresholds}

\subsection{Input, normalization, and scope}\label{input-normalization-and-scope}

Let \(a_n=A082161(n)\), and let \(z=a_1^{\operatorname{Ai}}=-2.338107410459767...\) denote the first negative zero of the Airy function. To avoid confusing this zero with the first sequence term, this report uses \(z\) throughout. The forward theorem provides a fixed \(\gamma>0\) satisfying

\[
\log a_n=\log\gamma+\log\Gamma(n+1)+n\log4+3z n^{1/3}+\log n
+B n^{-1/3}+D n^{-2/3}+L_3 n^{-1}+O(n^{-4/3}),
\tag{I1}
\]

where

\[
B=\frac{53z^2}{90},\qquad D=\frac{44z}{27},\qquad
L_3=\frac{141}{140}-\frac{1304z^3}{42525}.
\tag{I2}
\]

The error in (I1) is initially only an assertion at integer \(n\). All conclusions about an actual sequence below use only that assertion, not unproved differentiability of its remainder.

There are three distinct inversion objects:

\begin{enumerate}
\def\labelenumi{\arabic{enumi}.}
\tightlist
\item
  the inverse of an explicit smooth asymptotic model;
\item
  the inverse of a chosen interpolation of the integer sequence;
\item
  the discrete threshold \(N(X)=\min\{n:a_n\ge X\}\).
\end{enumerate}

They have different error statements. In particular, no real-valued asymptotic expansion, by itself, proves a single unqualified floor/ceiling rule at thresholds arbitrarily close to sequence values.

\subsection{The factorial contributes a crucial logarithm}\label{the-factorial-contributes-a-crucial-logarithm}

Stirling's formula gives

\[
\log\Gamma(n+1)=n\log n-n+\tfrac12\log n+
\tfrac12\log(2\pi)+\frac1{12n}+O(n^{-3}).
\]

Consequently, set

\[
A=3z,\quad p=\tfrac32,\quad C=\log(\gamma\sqrt{2\pi}),\quad
E=L_3+\tfrac1{12}=\frac{229}{210}-\frac{1304z^3}{42525},
\tag{I3}
\]

and define

\[
h(n)=n\log(4n/e)+A n^{1/3}+p\log n+C
+B n^{-1/3}+D n^{-2/3}+E n^{-1}.
\tag{I4}
\]

Then (I1) is \(\log a_n=h(n)+O(n^{-4/3})\). In particular, the coefficient of \(\log n\) is \textbf{3/2}, not 1. Numerically,

\[
\begin{aligned}
A&=-7.014322231379301\ldots,& B&=3.219306132565383\ldots,\\
D&=-3.810249113341843\ldots,& E&=1.482422558319535\ldots.
\end{aligned}
\]

An equally useful model retains the factorial exactly:

\[
h_\Gamma(n)=\log\gamma+\log\Gamma(n+1)+n\log4+A n^{1/3}+\log n
+B n^{-1/3}+D n^{-2/3}+L_3/n.
\tag{I5}
\]

Its difference from (I4) is \(-1/(360n^3)+O(n^{-5})\), below the precision needed here. Its inverse differs from the inverse of (I4) by \(O(n^{-3}/\log n)\).

For (I4),

\[
h'(n)=\log(4n)+\frac A3n^{-2/3}+\frac pn
-\frac B3n^{-4/3}-\frac{2D}3n^{-5/3}-E n^{-2}
\sim\log(4n)>0,
\tag{I6}
\]

and \(h''(n)\sim 1/n\). Thus \(h\) has a unique increasing large-argument inverse.

\subsection{Lambert W anchor and the first Airy correction}\label{lambert-w-anchor-and-the-first-airy-correction}

Put

\[
T=\log X,\qquad x=\frac{T}{W_0(4T/e)},\qquad
s=\log(4x)=1+W_0(4T/e),\qquad q=p\log x+C.
\tag{I7}
\]

Here \(W_0\) is the principal real Lambert W branch and \(X\) is sufficiently large. The identity

\[
x\log(4x/e)=T
\]

holds exactly. The leading correction needed to undo the negative Airy term is positive:

\[
h^{-1}(T)=x-\frac{A}{s}x^{1/3}-\frac{q}{s}
+O\!\left(\frac{x^{-1/3}}s\right).
\tag{I8}
\]

Equivalently, the displayed first terms are

\[
x-\frac{3z}{\log(4x)}x^{1/3}-\frac32
+\frac{\frac32\log4-\log(\gamma\sqrt{2\pi})}{\log(4x)}.
\tag{I9}
\]

This shows the hierarchy clearly: the Airy correction grows like \(x^{1/3}/\log x\), the factorial/prefactor combination gives a constant shift \(-3/2\), and the unknown amplitude first enters at order \(1/\log x\).

Equation (I8), including its stated error scale for suitable interpolations or the discrete bracketing interpretation below, already follows from the weaker core hypothesis

\[
a_n=\gamma n!4^n e^{3z n^{1/3}}n\bigl(1+O(n^{-1/3})\bigr).
\]

It does not require the all-orders construction.

An alternative anchor explicitly absorbs the unknown amplitude: set \(y=T-C\), \(x=y/W_0(4y/e)\), and \(q=p \log x\). Every formula below remains valid with this replacement. The unshifted version (I7) makes the dependence on \(\gamma\) easier to see.

\subsection{Explicit inverse through the third fractional order}\label{explicit-inverse-through-the-third-fractional-order}

Define

\[
\begin{aligned}
u={}&-\frac A s,\\
v={}&-\frac q s,\\
w={}&-\frac B s+\frac{A^2}{3s^2}-\frac{A^2}{2s^3},\\
r={}&-\frac D s+\frac{A(p+q/3)}{s^2}-\frac{Aq}{s^3},\\
k={}&-\frac E s+\frac{pq}{s^2}-\frac{AB+q^2/2}{s^3}
+\frac{A^3}{3s^4}-\frac{A^3}{2s^5}.
\end{aligned}
\tag{I10}
\]

Then the inverse of the explicit smooth model (I4) satisfies

\[
\boxed{\displaystyle
h^{-1}(T)=x+u x^{1/3}+v+w x^{-1/3}+r x^{-2/3}+k x^{-1}
+O\!\left(\frac{x^{-4/3}}{\log(4x)}\right).}
\tag{I11}
\]

The error order in (I11) is also the uncertainty imposed by the forward remainder in (I1). It is therefore the natural stopping point using the three displayed forward logarithmic coefficients.

\subsubsection{Derivation}\label{derivation}

Write \(t=x^{-1/3}\) and

\[
\delta=u/t+v+wt+rt^2+kt^3,\qquad n=x+\delta.
\]

The leading function \(f(n)=n \log(4n/e)\) has expansion

\[
f(x+\delta)-f(x)=s\delta+\frac{\delta^2}{2x}
-\frac{\delta^3}{6x^2}+\cdots.
\]

For the perturbations, use

\[
\begin{aligned}
A(x+\delta)^{1/3}&=A/t+(A/3)ut+(A/3)vt^2
+\bigl((A/3)w-(A/9)u^2\bigr)t^3+O(t^4),\\
p\log(x+\delta)+C&=q+pu t^2+pv t^3+O(t^4),\\
B(x+\delta)^{-1/3}&=Bt-(B/3)u t^3+O(t^4),\\
D(x+\delta)^{-2/3}&=Dt^2+O(t^4),\\
E(x+\delta)^{-1}&=Et^3+O(t^5).
\end{aligned}
\]

The required equations, in order, are

\[
\begin{aligned}
su+A&=0,\\
sv+q&=0,\\
sw+u^2/2+Au/3+B&=0,\\
sr+uv+Av/3+pu+D&=0,\\
sk+v^2/2+(u+A/3)w-u^3/6-Au^2/9+pv-Bu/3+E&=0.
\end{aligned}
\tag{I12}
\]

Solving (I12) gives (I10); the substantial cancellation in \(k\) is exact. SymPy verifies every coefficient from \(t^{-1}\) through \(t^3\) is identically zero with \(A,B,D,E,p,q,s\) left symbolic.

Although \(q\) depends on \(x\), this calculation is a Taylor expansion about the fixed anchor \(x\). Its contribution to derivatives is already included in \(p \log(x+\delta)\); one must not separately differentiate \(q\) within (I12).

For the remainder, \(q=O(s)\), \(u=O(1/s)\), \(v=O(1)\), and \(w,r,k=O(1/s)\). Hence \(\delta=O(x^{1/3}/s)\) and \(\delta/x=O(x^{-2/3}/s)\). All omitted terms in the residual \(h(x+\delta)-T\) are \(O(x^{-4/3})\). Since the local derivative is asymptotic to \(s\), the mean-value theorem yields (I11). No derivative bound on the unknown sequence remainder was used.

\subsection{Two Newton steps are a simpler practical inverse}\label{two-newton-steps-are-a-simpler-practical-inverse}

For computation it is usually preferable to avoid manually evaluating (I10). Starting from the Lambert anchor \(x_0=x\), take

\[
x_{j+1}=x_j-\frac{h(x_j)-T}{h'(x_j)},\qquad j=0,1.
\tag{I13}
\]

If \(n_*=h^{-1}(T)\), then

\[
\begin{aligned}
x_0-n_*&=O(x^{1/3}/s),\\
x_1-n_*&=O(x^{-1/3}/s^3),\\
x_2-n_*&=O(x^{-5/3}/s^7).
\end{aligned}
\tag{I14}
\]

Indeed, throughout the relevant interval, \(h'\) is comparable to \(s\) and \(h''=O(1/x)\), so one Newton step squares the error and divides it by \(xs\). The last bound is smaller than the forward uncertainty \(O(x^{-4/3}/s)\). Thus two Newton steps solve the available asymptotic inverse more accurately than its supplied coefficients justify.

The same procedure works using \(h_\Gamma\) from (I5), with

\[
h_\Gamma'(n)=\psi(n+1)+\log4+(A/3)n^{-2/3}+n^{-1}
-(B/3)n^{-4/3}-(2D/3)n^{-5/3}-L_3n^{-2},
\]

where \(\psi\) is the digamma function. Work with \(T=\log X\) rather than forming very large \(X\).

More generally, if a smooth model has forward logarithmic error \(O(n^{-\rho})\), its inverse uncertainty is \(O(x^{-\rho}/s)\). Additional Newton steps can make model inversion error smaller than any prescribed finite power. Starting from (I7), after \(j\) steps the bound is

\[
x_j-n_*=O\!\left(x^{\,1-2^{j+1}/3}s^{\,1-2^{j+1}}\right).
\]

\subsection{What ``smooth inverse'' means for an integer sequence}\label{what-smooth-inverse-means-for-an-integer-sequence}

The asymptotic assertion (I1) alone does not specify a canonical interpolation between integers. It is incorrect simply to differentiate its \(O(n^{-4/3})\) remainder as though a differentiable real-variable asymptotic theorem had been proved.

There is, however, a precise conditional statement:

\begin{itemize}
\item
  If an eventually increasing interpolation \(L(t)\) satisfies \(L(n)=\log a_n\) and \(L(t)=h(t)+O(t^{-4/3})\) uniformly for real \(t\), then

  \[
  L^{-1}(T)=h^{-1}(T)+O(x^{-4/3}/s).
  \tag{I15}
  \]

  The argument is value bracketing using monotonicity and the slope of \(h\); no derivative estimate for the remainder of \(L\) is necessary.
\item
  Such smooth increasing interpolations \textbf{exist} under (I1). For example, put \(\varepsilon_n=\log a_n-h(n)\) and choose a fixed smooth bump \(\phi\), supported in \((-1/3,1/3)\), with \(\phi(0)=1\). For large real \(t\), set

  \[
  L(t)=h(t)+\sum_n\varepsilon_n\phi(t-n).
  \]

  The supports are disjoint, the perturbation and its first derivative are both \(O(t^{-4/3})\), and therefore \(L'(t)>0\) eventually. This supplies a smooth interpolation meeting (I15), but it is an arbitrary choice, not an intrinsic continuous extension of the combinatorial sequence.
\item
  The common interpolation that is piecewise linear in \(\log a_n\) has a weaker generic comparison. Because \(h''(t)\sim 1/t\), its interpolation error relative to \(h\) is \(O(1/t)\), larger than the forward remainder. Its inverse satisfies only

  \[
  L_{\mathrm{linear}}^{-1}(T)=h^{-1}(T)+O(1/(x\log x)).
  \tag{I16}
  \]
\item
  An arbitrary monotone interpolation can differ by order 1 in its inverse. Integer samples alone do not force the sharper real-valued estimate.
\end{itemize}

Every increasing interpolation through the exact sequence nevertheless gives the same discrete threshold after taking a ceiling, with \(N(X)=\operatorname{ceil}(L^{-1}(\log X))\) for large \(X\). Thus interpolation distinctions do not obstruct the direct discrete bracketing below.

\subsection{Discrete threshold bracketing and the rounding barrier}\label{discrete-threshold-bracketing-and-the-rounding-barrier}

The forward hypothesis itself implies eventual strict monotonicity. In fact,

\[
\log a_{n+1}-\log a_n=\log(4n)+O(n^{-2/3})>0
\]

under (I1). Under just the weaker core hypothesis, the error can instead be written \(O(n^{-1/3})\), still sufficient. As \(a_n\) is unbounded, finitely many earlier terms do not affect \(N(X)\) for sufficiently large \(X\).

Suppose that an actual forward bound is available:

\[
|\log a_n-h(n)|\le K n^{-4/3}\qquad(n\ge n_0).
\tag{I17}
\]

Let \(r_*=h^{-1}(\log X)\) and \(K_*=\max(K,1)\). For all sufficiently large \(r_*\), one valid, deliberately conservative choice is

\[
\eta=\frac{8K_* r_*^{-4/3}}{\log(4r_*)}.
\]

Then

\[
\boxed{\quad
\lceil r_*-\eta\rceil\ \le\ N(X)\ \le\ \lceil r_*+\eta\rceil.
\quad}
\tag{I18}
\]

To verify (I18), test the integer immediately below \(\operatorname{ceil}(r_*-\eta)\) and the integer \(\operatorname{ceil}(r_*+\eta)\). They lie within a bounded distance of \(r_*\); there \(h'\) is at least half \(\log(4r_*)\), while the allowed forward error is at most twice \(K r_*^{-4/3}\). The displacement \(\eta\) leaves a strict margin on either side. Eventual monotonicity gives the result.

If the explicit approximation (I11) is used in place of \(r_*\), enlarge \(\eta\) by a bound for its error, which has the same order. If the two-step Newton approximation is used, its much smaller error from (I14) can be added instead.

Important qualifications:

\begin{enumerate}
\def\labelenumi{\arabic{enumi}.}
\tightlist
\item
  Once \(2 \eta<1\), (I18) leaves either one integer or two adjacent candidates.
\item
  If \(\operatorname{dist}(r_*,\mathbb Z)>\eta\), both ceilings agree and \(N(X)=\operatorname{ceil}(r_*)\).
\item
  If \(r_*\) is within \(\eta\) of an integer, the next omitted term or exact sequence evaluation can decide which adjacent candidate is correct. The supplied asymptotic does not always decide it.
\item
  One may write \(N(X)=\operatorname{ceil}(r_*+O(r_*^{-4/3}/\log r_*))\) only in the existential/bracketing sense of (I18). It does \textbf{not} mean \(N(X)=\operatorname{ceil}(r_*)+o(1)\) uniformly in \(X\).
\item
  An unquantified \(O\) bound proves the existence of (I18) with some constant and starting point. It does not yield a numerically certified enclosure for a particular finite \(X\). That requires effective forward constants, a certified amplitude, and numerical error control.
\end{enumerate}

Under the weaker core theorem, the corresponding bracket radius is \(O(r_*^{-1/3}/\log r_*)\), still tending to zero. Therefore the core theorem alone already localizes the threshold to at most two adjacent integers asymptotically.

\subsection{Unknown gamma and numerical amplitude estimates}\label{unknown-gamma-and-numerical-amplitude-estimates}

The constant \(\gamma\) is a fixed positive parameter, not a free constant that can be silently dropped after leading order. Changing \(\gamma\) to \(\widetilde{\gamma}\) shifts the smooth inverse, to first order, by

\[
-\frac{\log(\widetilde\gamma/\gamma)}{\log(4n)}.
\tag{I19}
\]

A bounded multiplicative amplitude uncertainty thus causes an inverse uncertainty \(O(1/\log n)\). This is \(o(1)\), but is much larger than the algebraic error in (I11), and can still change the discrete answer at near-integer thresholds.

If certified information gives

\[
|\log\gamma-\log\widetilde\gamma|\le\varepsilon,
\]

the forward error budget for a model using \(\widetilde{\gamma}\) becomes \(\varepsilon+K n^{-4/3}\). The threshold radius is then

\[
O\!\left(\frac{\varepsilon+r_*^{-4/3}}{\log(4r_*)}\right).
\]

An uncertified fitted decimal for \(\gamma\) does not supply such a bound. In particular, the supplied approximate value \(166.95063854245\) is suitable for illustrative computations, not for asserting rigorous threshold decisions.

\subsubsection{An amplitude-free locator and its existing provenance}\label{an-amplitude-free-locator-and-its-existing-provenance}

If \(\gamma\) is unknown, one can still use

\[
r_0=x-\frac{3z}{\log(4x)}x^{1/3}-\frac32,
\qquad x=\frac{\log X}{W_0(4\log X/e)}.
\]

The real inverse uncertainty is \(O(1/\log x)=o(1)\). In discrete form, for some fixed constant \(K\) and sufficiently large \(X\),

\[
\left\lceil r_0-\frac K{\log x}\right\rceil
\le N(X)\le
\left\lceil r_0+\frac K{\log x}\right\rceil.
\]

This coarse result already follows from the published Elvey Price--Fang--Wallner Θ theorem: its logarithm has an \(O(1)\) remainder, which is divided by the slope \(\sim \log x\) on inversion. It should therefore not be presented as a consequence unique to the amplitude proof. The amplitude convergence identifies the fixed \(\log \gamma\) term; the all-orders theorem provides the finer algebraic inverse corrections. As before, existence of an unspecified constant \(K\) is not a numerical finite-threshold certificate.

\section{Reproducible computation and numerical limitations}\label{reproducible-computation-and-numerical-limitations}

The accompanying source archive is relocatable and includes every script needed to generate finite-order coefficients, verify the algebraic identities used above, and reproduce the reported numerical evidence. It requires Python with SymPy, NumPy, SciPy, and mpmath. The exact versions used in the delivered replay are recorded in its environment report.

The main replay command is \texttt{python\ replay.py\ -\/-max-n\ 20000}. It verifies the SHA-256 manifest of its inputs before executing the checks, writes every result beneath \texttt{output/}, and reports any failure. The coefficient command \texttt{python\ coefficients.py\ -\/-order\ 8\ -\/-output\ output/coefficients.json} constructs all scalar and profile terms through the requested order. Raising the order invokes the same finite polynomial algorithm, with no lookup table or fitted coefficients.

The coefficient construction is a finite, exact procedure:

\begin{enumerate}
\def\labelenumi{\arabic{enumi}.}
\tightlist
\item
  Represent a profile by a polynomial pair \((P,Q)\) for \(PF+QF'\) and differentiate using \(F''=2(x+a)F\)
\item
  Expand the rational recurrence, the two shifted arguments, and the previous-time powers to the next required degree in \(\varepsilon\)
\item
  Collect the known forcing in the two components \(F\) and \(F'\)
\item
  Solve the triangular polynomial equation (F28), choose the scalar so that \(Q(0)=0\), and integrate for \(P\) with \(P(0)+Q'(0)=0\)
\item
  Verify the resulting residual coefficient, both boundary conditions, and the mod-three grading exactly
\item
  Match the logarithmic one-step ratio and expand the endpoint factor to obtain forward logarithmic and relative coefficients
\end{enumerate}

The archive also checks the original triangular recurrence against the transformed recurrence with exact rational arithmetic, verifies the squared gauge identities, verifies the finite frozen quasimode, and checks the inverse cancellations with symbolic parameters. Its numerical tests use the full finite support of the forward recurrence, rather than an unquantified height cutoff.

At \(n=20000\), the normalized raw ratio is \(187.01754328008607\). Dividing it by the first three relative corrections gives \(166.95063854245257\); removing the first three logarithmic corrections instead gives \(166.95152887373968\). These values are consistent with the empirical amplitude \(166.95208957\) reported by Elvey Price, Fang, and Wallner. They are not certified enclosures, and the last displayed digits should not be interpreted as established digits of \(\gamma\).

The exact symbolic checks establish the displayed cancellation identities. Floating-point checks establish only reproducibility of numerical evidence. Neither kind of computation replaces the analytic compactness, stability, or endpoint arguments above. The theorem's constant is the convergent expression (F25), not a fitted decimal. In particular, certified finite-threshold decisions require effective forward constants, a certified amplitude interval, and controlled numerical evaluation.

\section{Further research questions}\label{further-research-questions}

\textbf{Certified amplitude evaluation.} Can the convergent spectral-product and projection representation be supplemented with explicit tail constants to give an efficient interval enclosure of \(\gamma\)? The present proof gives existence and algebraic rates, but its compactness-based spectral gap has not been made numerically effective. A constructive lower gap, certified Airy residuals, and explicit endpoint bounds would make this a concrete certification problem.

\textbf{Compacted binary trees.} The corresponding compacted-tree recurrence has a subtraction term. Its established leading scale has the same stretched exponential and a different power of \(n\). Extending the present positive-kernel smoothing argument requires an appropriate positive lift or a different stability estimate. The relaxed-tree proof does not automatically establish the compacted amplitude.

\textbf{Higher arity.} The exact gauge and parity decomposition exploit a nearest-neighbor binary recurrence. For relaxed \(k\)-ary structures, one can ask whether a comparable slowly varying transfer operator admits a coercive continuum limit and summable ground-mode drift. The specific gauge, gap, and boundary exponent would need to be derived anew.

\textbf{Beyond algebraic orders.} The construction proves a Poincaré expansion at every fixed order. It gives no optimal-truncation theorem, no Gevrey estimate on the coefficient growth, and no exponentially small correction. Those questions require quantitative control uniform in the expansion order, together with an analysis of higher Airy modes and possible additional asymptotic sectors.

\textbf{Intrinsic interpolation and inverse certification.} The combinatorial sequence has no canonical real interpolation supplied by its recurrence. The threshold brackets in §13 avoid assuming one. It would be useful to combine interval evaluation of nearby exact terms with the asymptotic locator, yielding a certified algorithm that resolves the occasional two-candidate ambiguity.

\section{References}\label{references}

\begin{enumerate}
\def\labelenumi{\arabic{enumi}.}
\tightlist
\item
  Andrew Elvey Price, Wenjie Fang, and Michael Wallner. \emph{Compacted binary trees admit a stretched exponential}. Journal of Combinatorial Theory, Series A 177 (2021), article 105306. \href{https://arxiv.org/pdf/1908.11181}{arXiv:1908.11181}. The counting recurrence, the proved Θ estimate, and the conjectural amplitude and all-orders form used here originate in this paper.
\item
  The On-Line Encyclopedia of Integer Sequences. \href{https://oeis.org/A082161}{A082161}. Sequence identification and combinatorial interpretations.
\item
  NIST Digital Library of Mathematical Functions. \href{https://dlmf.nist.gov/9}{Chapter 9, Airy and Related Functions}. Standard Airy differential equation, zeros, and decay properties.
\end{enumerate}

The proposed convergence argument and its scope are distinguished from the existing result in §1. No assertion of an exhaustive global priority search or of peer-reviewed acceptance is made.

\end{document}
